\chapter[Fundamentação Teórica]{Fundamentação Teórica}

Este capítulo apresenta os conceitos necessários para compreender a análise e a predição de trajetórias de espermatozoides em vídeos microscópicos. A organização acompanha o fluxo dos dados na pipeline: o domínio de aplicação define as medidas de interesse; o processamento de imagens permite localizar as células; o rastreamento associa observações ao longo do tempo; o fluxo óptico descreve movimento aparente; e a predição estima posições ainda não observadas. Ao final, são apresentados os conjuntos de dados e as métricas que permitem avaliar cada módulo separadamente.

\section{Análise seminal assistida por computador}

A análise seminal examina características do sêmen e dos espermatozoides, entre elas concentração, contagem total, morfologia, vitalidade e motilidade. A sexta edição do manual da Organização Mundial da Saúde sistematiza procedimentos de laboratório e ressalta a importância do controle de qualidade para produzir resultados comparáveis \cite{who2021manual}. Na avaliação básica, a motilidade pode ser classificada como progressiva, não progressiva ou imóvel.

Sistemas \ac{CASA} combinam microscopia, aquisição digital de vídeo e algoritmos para tornar essa avaliação quantitativa. Em vez de resumir a amostra somente pela inspeção visual, esses sistemas localizam cabeças de espermatozoides, associam observações da mesma célula e calculam grandezas cinemáticas a partir das coordenadas registradas no tempo. A automação, entretanto, não elimina a necessidade de padronização: taxa de quadros, ampliação, profundidade da câmara, contraste e regras de classificação influenciam diretamente os valores obtidos.

Considere uma trajetória bidimensional formada pelas posições $\mathbf{p}_t=(x_t,y_t)$, observadas em intervalos constantes $\Delta t$. Para uma trajetória com $T$ posições, a velocidade curvilínea, \ac{VCL}, corresponde ao comprimento acumulado do caminho dividido pelo tempo, enquanto a velocidade em linha reta, \ac{VSL}, considera apenas a distância entre a primeira e a última posição:

\begin{align}
	\operatorname{VCL} &=
	\frac{\sum_{t=1}^{T-1}\lVert \mathbf{p}_{t+1}-\mathbf{p}_t\rVert_2}
	{(T-1)\Delta t},\\
	\operatorname{VSL} &=
	\frac{\lVert \mathbf{p}_{T}-\mathbf{p}_1\rVert_2}
	{(T-1)\Delta t}.
\end{align}

A velocidade do percurso médio, \ac{VAP}, é calculada ao longo de uma versão suavizada da trajetória. A linearidade, \ac{LIN}, compara o deslocamento líquido com o caminho curvilíneo e pode ser expressa por $\operatorname{LIN}=100\,\operatorname{VSL}/\operatorname{VCL}$ \cite{who2021manual}. Para que as velocidades sejam apresentadas em $\mu$m/s, é necessário conhecer a escala espacial e a taxa de quadros. Sem essa calibração, as medidas devem permanecer em pixels por unidade de tempo.

\section{Processamento de imagens e detecção}

Uma imagem digital pode ser representada por uma função discreta de intensidade ou por múltiplos canais de cor. Em vídeos microscópicos, o processamento inicial busca aumentar a separação entre as cabeças dos espermatozoides e o fundo. Conversão para tons de cinza, redução de ruído, normalização de contraste, limiarização e operações morfológicas são transformações comuns nessa etapa \cite{szeliski2022computer}.

A segmentação por limiar divide os pixels de acordo com sua intensidade. O método de Otsu seleciona automaticamente um limiar que maximiza a separação estatística entre classes do histograma \cite{otsu1979threshold}. Após a binarização, erosão e dilatação podem remover componentes pequenos, preencher falhas ou separar estruturas segundo o elemento estruturante. Quando objetos próximos formam uma única região, a transformação \textit{watershed} oferece uma alternativa baseada em marcadores. Esses métodos são rápidos e interpretáveis, mas seus parâmetros são sensíveis a mudanças de contraste, iluminação, densidade celular e presença de detritos.

Na detecção de objetos, a saída contém, para cada instância, uma classe, uma caixa delimitadora e uma confiança. A família YOLO consolidou detectores de um estágio que realizam localização e classificação em uma única rede neural \cite{redmon2016yolo}. O VISEM-Tracking utiliza coordenadas normalizadas em formato compatível com YOLO e oferece classes para espermatozoides, aglomerados e células pequenas ou com cabeça reduzida \cite{thambawita2023visemtracking}.

A detecção de espermatozoides é um caso de detecção de objetos pequenos: poucos pixels descrevem cada cabeça, a textura individual é limitada e o borramento de movimento reduz a nitidez. Por isso, abordagens recentes utilizam resoluções maiores, camadas dedicadas a objetos pequenos e mecanismos de atenção. Zhang et al. \cite{zhang2024spermyolo}, por exemplo, adaptam o YOLOv8 e um rastreador para o domínio espermático. A avaliação precisa considerar precisão, revocação, qualidade de localização e tempo de execução, pois os erros do detector são propagados para o rastreamento.

\section{Rastreamento multiobjeto}

O \ac{MOT} consiste em estimar simultaneamente o estado de vários objetos e manter um identificador consistente para cada um ao longo do vídeo. Na estratégia \textit{tracking-by-detection}, cada quadro é processado por um detector e as novas observações são associadas às trajetórias existentes. O estado pode incluir posição, tamanho da caixa e velocidade. Uma trajetória é iniciada quando uma detecção persiste, atualizada enquanto recebe observações compatíveis e encerrada após permanecer ausente por um número definido de quadros.

Dois fundamentos desse processo são a estimação de estado e a associação de dados. O Filtro de Kalman representa a evolução do estado por um modelo linear com incerteza: primeiro prediz o estado no novo instante e depois corrige a previsão com a medição observada, ponderando as covariâncias dos ruídos de processo e de medição \cite{kalman1960filter}. Para associar trajetórias e detecções, constrói-se uma matriz de custos com distâncias entre centroides, sobreposição entre caixas ou outra medida de similaridade. O algoritmo Húngaro fornece uma atribuição global de custo mínimo sob correspondências um a um.

O SORT combina Filtro de Kalman, sobreposição entre caixas e associação Húngara em um rastreador simples e rápido \cite{bewley2016sort}. O ByteTrack, por sua vez, realiza associações sucessivas com detecções de alta e baixa confiança, buscando recuperar objetos que seriam descartados por um limiar único \cite{zhang2022bytetrack}. Esses métodos mostram a evolução de modelos baseados principalmente em movimento para estratégias que combinam localização, confiança e, em alguns casos, aparência.

No domínio espermático, a associação é dificultada pela alta densidade, pelas trajetórias rápidas e curvas e pela semelhança visual entre indivíduos. Descritores de reidentificação podem ser pouco discriminativos quando as cabeças ocupam poucos pixels; por outro lado, utilizar somente distância espacial falha em cruzamentos e mudanças bruscas de direção. As identidades manuais do VISEM-Tracking permitem medir trocas de identidade, fragmentações e recuperação após ausências \cite{thambawita2023visemtracking}.

\section{Fluxo óptico e movimento aparente}

Fluxo óptico é um campo vetorial que descreve o deslocamento aparente de padrões de intensidade entre quadros. Para cada posição $(x,y)$, procura-se um vetor $\mathbf{v}(x,y)=(u,v)$ que relacione a imagem no instante atual à imagem seguinte. Sob a hipótese de constância de brilho,

\begin{equation}
	I(x,y,t)=I(x+u\Delta t,y+v\Delta t,t+\Delta t),
\end{equation}

uma aproximação de primeira ordem produz a equação de restrição

\begin{equation}
	I_xu+I_yv+I_t=0.
\end{equation}

Há uma equação para duas incógnitas, portanto são necessárias restrições adicionais. Métodos locais, como Lucas--Kanade, supõem movimento aproximadamente constante em uma vizinhança. Formulações globais, como Horn--Schunck, introduzem suavidade no campo. Farnebäck aproxima vizinhanças por polinômios e estima fluxo denso em múltiplas escalas, constituindo uma linha de base prática para vídeos \cite{farneback2003twoframe,szeliski2022computer}.

Abordagens aprendidas estimam correspondências a partir de bases com deslocamento conhecido. O RAFT constrói volumes de correlação entre características dos dois quadros e atualiza iterativamente o campo por uma unidade recorrente \cite{teed2020raft}. Embora apresente bons resultados em conjuntos gerais de fluxo óptico, sua transferência para microscopia precisa ser avaliada, pois escala, textura, contraste e formação da imagem diferem das cenas de treinamento.

Neste trabalho, o fluxo possui dois usos potenciais: compensar movimento global da imagem e produzir características locais para a predição. É necessário distinguir esse campo da velocidade física do fluido. O fluxo óptico é medido em coordenadas de imagem e pode misturar deslocamento das células, do meio visível, de artefatos e da câmera. Como os dados do VISEM e as anotações do VISEM-Tracking não fornecem uma medição física de referência do escoamento, o resultado será denominado campo de movimento aparente.

\section{Predição de trajetórias}

Predição de trajetória consiste em estimar posições futuras a partir de uma sequência observada. Se $\mathbf{p}_{t-p+1:t}$ representa uma janela de $p$ posições, o objetivo é produzir $\widehat{\mathbf{p}}_{t+1:t+q}$ para um horizonte de $q$ quadros. A tarefa difere do rastreamento: o rastreador estima o estado atual com medições já disponíveis, enquanto o preditor é avaliado em posições ocultadas durante a inferência.

Modelos cinemáticos simples são referências importantes. A extrapolação de velocidade constante usa o deslocamento recente para projetar a posição, e o Filtro de Kalman acrescenta incertezas de processo e medição \cite{kalman1960filter}. Esses métodos são interpretáveis e exigem poucos dados, mas tendem a degradar quando a célula altera rapidamente sua direção ou velocidade.

Redes recorrentes representam dependências em sequências. A \ac{LSTM} utiliza células de memória e portas de controle para preservar informação por intervalos mais longos \cite{hochreiter1997lstm}. Em predição de trajetórias, posições ou deslocamentos sucessivos são codificados em um estado oculto, a partir do qual são geradas coordenadas futuras. O Social LSTM mostrou como essa formulação pode ser aplicada a trajetórias de agentes \cite{alahi2016sociallstm}; no domínio deste trabalho, Noy et al. \cite{noy2023location} aplicaram LSTM à localização futura de espermatozoides.

A hipótese computacional desta pesquisa considera que o histórico geométrico pode ser complementado por movimento aparente. Um modelo sem fluxo pode receber $(x,y,\Delta x,\Delta y)$, enquanto o modelo com fluxo acrescenta $(u,v,\lVert\mathbf{v}\rVert,\theta)$ amostrados na vizinhança da célula. A comparação deve manter arquitetura, histórico, horizonte e divisão dos dados controlados, alterando somente as características de fluxo.

\section{Conjuntos de dados}

O VISEM contém 85 vídeos microscópicos, um por participante, associados a resultados de análise seminal e atributos biológicos, mas sua publicação original não fornece identificadores de rastreamento quadro a quadro \cite{haugen2019visem}. O VISEM-Tracking acrescenta caixas delimitadoras, classes e identificadores persistentes verificados por especialistas a trechos iniciais de 30 segundos de 20 desses vídeos, totalizando 29.196 quadros \cite{thambawita2023visemtracking}. Trata-se de uma extensão anotada da mesma coleção de vídeos. Neste trabalho, os trechos anotados serão utilizados no desenvolvimento e na avaliação quantitativa de detecção, rastreamento e predição. Os demais 65 vídeos serão destinados à aplicação posterior e à análise qualitativa, sem atribuir acurácia às trajetórias na ausência de anotações manuais.

O protocolo de aquisição descreve amostras em platina aquecida a $37\,^{\circ}\mathrm{C}$, examinadas em microscópio Olympus CX31 com aumento de $400\times$ e registradas por câmera UEye UI-2210C acoplada ao microscópio \cite{haugen2019visem,thambawita2023visemtracking}. Essas publicações não documentam tubo, bombeamento ou contracorrente induzida durante a filmagem. O artigo original relata que, para alguns participantes, foram realizadas duas gravações devido à deriva na primeira amostra registrada, mantendo-se um vídeo por participante na coleção \cite{haugen2019visem}. Esse relato não estabelece a origem física da deriva nem comprova um escoamento controlado; a ausência de tal descrição também não demonstra que o líquido esteja imóvel.

As divisões de treinamento, validação e teste serão feitas por vídeo de origem, mantendo os trechos do mesmo participante na mesma divisão, nunca por quadros aleatórios. Clipes derivados não serão tratados como amostras independentes. Essa precaução impede que imagens quase idênticas da mesma sequência apareçam simultaneamente no treinamento e no teste.

\section{Métricas de avaliação}

Na detecção, precisão e revocação são definidas a partir de verdadeiros positivos (VP), falsos positivos (FP) e falsos negativos (FN):

\begin{equation}
	\text{Precisão}=\frac{\mathrm{VP}}{\mathrm{VP}+\mathrm{FP}},
	\qquad
	\text{Revocação}=\frac{\mathrm{VP}}{\mathrm{VP}+\mathrm{FN}}.
\end{equation}

Uma detecção é considerada correspondente ao gabarito quando a \ac{IoU} ultrapassa um limiar. A área sob a curva de precisão--revocação gera a precisão média, e a média entre classes ou limiares gera a \ac{mAP}. Como poucos pixels de deslocamento afetam significativamente a IoU de objetos pequenos, também é útil analisar erro de centroide e resultados por classe.

No rastreamento, a \ac{MOTA} combina falsos positivos, falsos negativos e trocas de identidade. A \ac{IDF1} mede a preservação de identidade por meio da média harmônica entre precisão e revocação no espaço das associações. A \ac{HOTA} foi proposta para equilibrar qualidade de detecção, associação e localização \cite{luiten2021hota}. Neste trabalho, HOTA será a métrica principal de rastreamento, acompanhada por MOTA, IDF1, trocas de identidade, fragmentações e tempo por quadro.

Na predição, o \ac{ADE} resume o erro ao longo de todo o horizonte, enquanto o \ac{FDE} considera a última posição prevista. Para $N$ trajetórias e horizonte de $q$ quadros:

\begin{align}
	\operatorname{ADE} &=
	\frac{1}{Nq}\sum_{i=1}^{N}\sum_{k=1}^{q}
	\left\lVert\widehat{\mathbf{p}}_{i,t+k}-\mathbf{p}_{i,t+k}\right\rVert_2,\\
	\operatorname{FDE} &=
	\frac{1}{N}\sum_{i=1}^{N}
	\left\lVert\widehat{\mathbf{p}}_{i,t+q}-\mathbf{p}_{i,t+q}\right\rVert_2.
\end{align}

Essas medidas devem ser apresentadas no mesmo sistema de coordenadas e, quando houver calibração, também em unidade física \cite{alahi2016sociallstm}. Médias serão acompanhadas por dispersão, intervalos de confiança e análise estratificada por horizonte, densidade celular e duração da trajetória.
